\section{Exact Kummer images and diagonal Fermat fibres}\label{sec:kummer}
The Galois actions in the manuscript can be made explicit by choosing
compatible roots and keeping the cyclotomic multiplier in the affine
group. The choice of origin changes the translation coordinate, so
we specify it throughout.

Write $\widehat\Z=\varprojlim_N\Z/N$ and put
\[
 \operatorname{Aff}_1(\widehat\Z)
 =\widehat\Z\rtimes\widehat\Z^\times,
 \qquad (a,b)(a',b')=(aa',b+ab'),
\]
where $(a,b)$ acts by $x\mapsto ax+b$. Take
$\zeta_N=e^{2\pi i/N}$ and the positive real roots
$\theta_N=2^{1/N}$, compatibly under divisibility. An automorphism
$\sigma\in\operatorname{Gal}(\overline\Q/\Q)$ then has coordinates
\[
 \sigma(\zeta_N)=\zeta_N^a,\qquad
 \sigma(\theta_N)=\zeta_N^b\theta_N.
\]
For odd $a$, define two characters by
\[
\begin{array}{c|rrrr}
 a\bmod8&1&3&5&7\\\hline
 \chi_2(a)&1&-1&-1&1\\
 \chi_{-2}(a)&1&1&-1&-1.
\end{array}
\]
The exponent $(-1)^b$ depends on $b\bmod2$.

\subsection{The radicals of two and minus one half}
\begin{proof}[Proof of the rank-one assertions of Theorem~\ref{thm:kummer-main}]
The identity $\sqrt2=\zeta_8+\zeta_8^{-1}$ gives the necessary
relation $(-1)^b=\chi_2(a)$. To prove sufficiency, put
$L=\Q(\mu_\infty)$. We first determine precisely the radical
relations of $2$ in $L$.

If $2^j$ is an $N$th power in $L$, then $2^{j/N}\in L$, since
$L$ contains all roots of unity. Reduce $j/N$ to $u/d$ with
$\gcd(u,d)=1$. B\'ezout's identity gives
$\Q(2^{u/d})=\Q(2^{1/d})$. The polynomial $X^d-2$ is Eisenstein
at $2$. For $d>2$, its positive real root generates a real field
which omits a nonreal conjugate, so that field is not normal over
$\Q$. Every finite subfield of the abelian extension $L/\Q$ is
normal. Hence $d$ is $1$ or $2$. Conversely $1$ and $\sqrt2$
belong to $L$.

Here is the finite Kummer argument in the form needed. If a field
$F$ contains $\mu_N$, adjoining $N$th roots of elements $u_1,\ldots,u_r$
gives a Galois group embedded in $(\Z/N)^r$. Under the usual
character pairing, its annihilator consists exactly of those
$(j_1,\ldots,j_r)$ for which $\prod_i u_i^{j_i}\in F^{\times N}$:
the corresponding product of the chosen radicals is fixed exactly
in that case. The perfect pairing on the finite group then determines
the Galois image from these relations. In rank one, the preceding
paragraph shows that the translation image over $L$ is the image of
$2\Z/N$ for every $N$. Taking inverse limits gives $2\widehat\Z$.

The cyclotomic projection of the full Galois group is
$\widehat\Z^\times$. Each fibre over a multiplier is therefore a
coset of $2\widehat\Z$, and the necessary parity relation selects
that entire coset. This proves that the image is exactly
\[
 H_2=\{(a,b):(-1)^b=\chi_2(a)\}.
\]
The compatible roots of $1/2$ are $\theta_N^{-1}$, and their
translation coordinate is $-b$. Their image is the same subgroup.

At finite level $N$, if $8\mid N$ both the character and the parity
are determined by the reduced coordinates, giving index two.
If $8\nmid N$, every pair in
$(\Z/N)^\times\ltimes\Z/N$ lifts to $H_2$. For odd $N$ one can
choose the parity of $b$ freely. For $v_2(N)=1$ or $2$, the unused
residue of $a$ modulo $8$ can be chosen to give either value of
$\chi_2$; the Chinese remainder theorem imposes the remaining odd
congruences. Thus, putting $d_N=2$ when $8\mid N$ and $d_N=1$
otherwise, the finite image has order $N\varphi(N)/d_N$.

For $-1/2$, choose $r_N=\zeta_{2N}\theta_N^{-1}$.
These roots are compatible and satisfy $r_N^N=-1/2$.
The action has translation coordinate
\[
 c=t(a)-b,\qquad t(a)=(a-1)/2\in\widehat\Z.
\]
Division is well defined: $a-1$ belongs to $2\widehat\Z$ and
multiplication by $2$ is injective on $\widehat\Z$.
The identity $t(aa')=t(a)+at(a')$ shows that
$(a,b)\mapsto(a,t(a)-b)$ is an automorphism of the affine group.
Since $(-1)^{t(a)}\chi_2(a)=\chi_{-2}(a)$, its image is
\[
 H_{-2}=\{(a,c):(-1)^c=\chi_{-2}(a)\}.
\]
The same lifting argument gives the same finite orders and indices.
\end{proof}

\subsection{The three Fermat fibres}
Consider the Fermat curve
\[
 u_0^N+u_1^N+u_\infty^N=0
\]
and its map to the projective line with coordinate
$z=-u_0^N/u_\infty^N$. Over each of $z=2,-1,1/2$, choose the
coordinate whose $N$th power is the exceptional one and divide
all projective coordinates by it. The other two coordinates then
have $N$th power $-1/2$. Consequently each fibre is a product of
two copies of the same $-1/2$ Kummer torsor.

\begin{proposition}\label{prop:fermat-fibre}
With equal origins $r_N$ in both factors, the Galois image acts by
\[
 (x,y)\longmapsto(ax+c,ay+c),\qquad (a,c)\in H_{-2,N}.
\]
Its index in the scalar affine group
$(\Z/N)^2\rtimes(\Z/N)^\times$ is $Nd_N$, and its index in
$(\Z/N)^2\rtimes\operatorname{GL}_2(\Z/N)$ is
\[
 d_N N^4\prod_{\ell\mid N}(1-\ell^{-2}).
\]
Its adelic index in either ambient group is infinite. If the first
origin is instead $\zeta_N^{-1}r_N$, the first translation coordinate
is $c-a+1$ and the second remains $c$.
\end{proposition}
\begin{proof}
Equal origins give identical translation cocycles, proving the
diagonal formula. The representation on a fibre is faithful on
the rank-one image, whose order is $N\varphi(N)/d_N$.
The scalar ambient group has order $N^2\varphi(N)$, giving the
first index. For a prime power $\ell^k$, reduction from
$\operatorname{GL}_2(\Z/\ell^k)$ to
$\operatorname{GL}_2(\F_\ell)$ has kernel of order
$\ell^{4(k-1)}$, while
$|\operatorname{GL}_2(\F_\ell)|=(\ell^2-1)(\ell^2-\ell)$.
Multiplying over primes yields
\[
 |\operatorname{GL}_2(\Z/N)|
 =N^4\prod_{\ell\mid N}(1-\ell^{-1})(1-\ell^{-2}).
\]
Dividing the full affine order by the image order proves the second
formula. The unbounded indices $Nd_N$ force infinite adelic index.
For the changed origin, $x$ is replaced by $x+1$; writing $x'=x+1$, the
affine action becomes $x'\mapsto a(x'-1)+c+1$, hence the
stated coordinate.
\end{proof}
At levels $N=1,2,4,8$, the image orders are respectively $1,2,8,16$
and the scalar fibre indices are $1,2,4,16$. The first departure from
the unrestricted rank-one affine group is therefore visible at level
$8$, whereas the diagonal restriction is already visible at level $2$.

\subsection{A uniform bound for independent radicals}
The general openness assertion is classical Kummer theory. Schinzel's
abelian-binomial theorem supplies the bounded radical obstruction;
Perucca and Sgobba state it explicitly in
\cite[Theorem~15, p.~7]{PeruccaSgobba} and prove uniform degree bounds
in their Theorems~1 and~13, pp.~1 and~6--8. More recently, Perucca
\cite[Theorem~26, pp.~13--14]{Perucca2026} describes the degree-to-index
ratio for all radical levels of a finitely generated subgroup over a
number field from finitely many levels and explicit cyclotomic factors.
That result concerns the growth across all levels. We give the elementary
cohomological argument needed for the explicit saturation bound here.

Let $K$ be a number field, and let $\xi_1,\ldots,\xi_r\in K^\times$
be independent modulo roots of unity. Choose a finite set of finite
places $S$ outside which they are units. The group
\[
 U_S=\mathcal O_{K,S}^{\times}/\mu_K
\]
is a finite free abelian group by the $S$-unit theorem
\cite[Theorem~5.11, p.~90]{MilneANT}. Let $M\subset U_S$ be the
rank-$r$ subgroup spanned by their classes, and put
\[
 M^{\mathrm{sat}}=(M\otimes\Q)\cap U_S.
\]
Choose a positive integer $e$ annihilating the finite quotient
$M^{\mathrm{sat}}/M$, and put $w=|\mu_K|$ and $D=ew$.

\begin{proposition}\label{prop:general-kummer-bound}
Over $L=K(\mu_\infty)$, the inverse-limit Kummer translation image
contains $D\widehat\Z^r$. Over $K$, the scalar affine image has
index at most
\[
 [\widehat\Z^\times:\chi(G_K)]D^r
 =[K\cap\Q(\mu_\infty):\Q]D^r.
\]
It is therefore open in
$\widehat\Z^r\rtimes\widehat\Z^\times$.
For $r\ge2$ it has infinite index in
$\widehat\Z^r\rtimes\operatorname{GL}_r(\widehat\Z)$.
\end{proposition}
\begin{proof}
Fix $N$ and write $G=\operatorname{Gal}(L/K)$, an abelian group.
For a multiplicative cocycle $z:G\to\mu_N$, commutativity gives
\[
 z(t)^{\chi(s)-1}=z(s)^{\chi(t)-1}.
\]
For fixed $s$, the right side is the coboundary of $z(s)$.
Hence $\chi(s)-1$ annihilates $H^1(G,\mu_N)$ for every $s$.
The ideal of $\Z/N$ generated by these differences has annihilator
exactly $\mu_N(K)$ inside the cyclic module $\mu_N$.
Its generator is therefore $d=|\mu_N(K)|=\gcd(N,w)$.
It follows that $d$, and thus $w$, annihilates $H^1(G,\mu_N)$.
All cocycles in question are continuous; the calculation takes place
in their finite image and needs no finiteness assumption on $G$.

Suppose $u=\prod_j\xi_j^{a_j}=\alpha^N$ with $\alpha\in L$.
The cocycle $z(s)=s(\alpha)/\alpha$ has $w$th power a coboundary,
so for some $\zeta\in\mu_N$ the element
$y=\alpha^w/\zeta$ lies in $K$. Thus
\[
 y^N=u^w.
\]
Valuations outside $S$ show that $y$ is an $S$-unit, and in $U_S$
we have $N[y]=w\sum_j a_j[\xi_j]$. In particular
$[y]\in M^{\mathrm{sat}}$. Multiplication by $e$ puts $[y]$ in
$M$, whose basis is the independent list $[\xi_j]$. Comparing
coefficients in
\[
 N e[y]=D\sum_j a_j[\xi_j]
\]
gives $N\mid Da_j$ for every $j$.

By the finite Kummer pairing proved above, the annihilator of the
translation image in $(\Z/N)^r$ is contained in the kernel of
multiplication by $D$. Taking annihilators reverses this inclusion,
so the image contains $D(\Z/N)^r$. A closed subgroup of
$\widehat\Z^r$ is the intersection of the inverse images of its
finite quotient images. Consequently the inverse-limit image
contains $D\widehat\Z^r$.

The projection to multipliers is exactly $\chi(G_K)$, of index
$[K\cap\Q(\mu_\infty):\Q]$ in $\widehat\Z^\times$.
The translation subgroup has index at most $D^r$. Multiplying the
indices proves the asserted bound. When $r\ge2$, all linear parts
are scalar matrices. Modulo $\ell^k$, scalar matrices have order
$\varphi(\ell^k)$ while
$|\operatorname{GL}_r(\Z/\ell^k)|$
grows as $\ell^{r^2k}$ up to a factor independent of $k$.
Their index is unbounded, proving the last assertion. For $r=0$
the statement reduces to the cyclotomic image alone.
\end{proof}
